\begin{env}[4.3.3]
\EGAReaderTarget{canon-ref.0.4.3.3-fr}{0.4.3.3-fr}
Sean $\mathcal{G}_1$, $\mathcal{G}_2$ dos $\mathcal{B}$-Módulos. De la
definición de las imágenes inversas de haces de grupos abelianos (3.7.1)
se deduce de inmediato un homomorfismo canónico
\[
  \psi^*(\mathcal{G}_1)\otimes_{\psi^*(\mathcal{B})}\psi^*(\mathcal{G}_2)
  \longrightarrow
  \psi^*(\mathcal{G}_1\otimes_{\mathcal{B}}\mathcal{G}_2)
\]
de $\psi^*(\mathcal{B})$-Módulos. Como el tallo en un punto de un producto
tensorial de haces es el producto tensorial de los tallos en ese punto
(G, II, 2.8), se deduce de (3.7.2) que el homomorfismo anterior es, de
hecho, un \emph{isomorfismo}. Al tomar el producto tensorial con
$\mathcal{A}$, se obtiene un \emph{isomorfismo canónico} funtorial
\[
  \Psi^*(\mathcal{G}_1)\otimes_{\mathcal{A}}\Psi^*(\mathcal{G}_2)
  \xrightarrow{\sim}
  \Psi^*(\mathcal{G}_1\otimes_{\mathcal{B}}\mathcal{G}_2).
  \tag{4.3.3.1}\label{0.4.3.3.1-fr}
\]
\end{env}
